% Generated draft: requires mathematical and source review.
\documentclass[11pt]{article}
\usepackage[margin=25mm]{geometry}
\usepackage{fontspec}
\setmainfont{Latin Modern Roman}
\usepackage{amsmath,amssymb,mathtools,unicode-math}
\setmathfont{Latin Modern Math}
\usepackage{enumitem,needspace,xurl,hyperref}
\setlength{\parindent}{0pt}
\setlength{\parskip}{5pt}
\setlength{\emergencystretch}{3em}
\newcommand{\N}{\mathbb{N}}
\newcommand{\Z}{\mathbb{Z}}
\newcommand{\Q}{\mathbb{Q}}
\newcommand{\R}{\mathbb{R}}
\newcommand{\C}{\mathbb{C}}
\newcommand{\F}{\mathbb{F}}
\newcommand{\A}{\mathbb{A}}
\newcommand{\PP}{\mathbb{P}}
\newcommand{\E}{\mathbb{E}}
\newcommand{\eps}{\varepsilon}
\newcommand{\dd}{\,\mathrm{d}}
\begin{document}
\section*{UnsolvedMath ID: 800340}
\textbf{Source problem code:} catalogue entry
\subsection{Definitions and mathematical statement}
For a region \(R\), let \(\mathcal{T}(R)\) denote its set of admissible tilings under the intended tile shapes and covering rules, and write \(T(R)=|\mathcal{T}(R)|\) when this set is finite. Call a nonnegative integer square-ish if it belongs to
\[
\{a^2:a\in\mathbb{Z}_{\geq 0}\}\ \cup\ \{2a^2:a\in\mathbb{Z}_{\geq 0}\}.
\]
The supplied excerpt tentatively suggests a two-parameter assertion: taking an intrusion of length \(m\) in a \(2n\)-by-\(2n\) square always produces a region whose number of tilings is square-ish. If \(R_{n,m}\) denotes the intended region and \(D\) denotes the intended set of admissible parameter pairs, this proposed assertion has the schematic form
\[
\forall (n,m)\in D\quad \exists a\in\mathbb{Z}_{\geq 0}\quad
T(R_{n,m})=a^2\ \text{or}\ T(R_{n,m})=2a^2.
\]
Neither \(R_{n,m}\) nor \(D\) is defined by the supplied record. In particular, the excerpt does not specify the geometry, placement, or meaning of the length of an intrusion; the tile shapes and tiling rules; or the admissible parameter ranges and boundary cases. Thus this display records the proposed numerical property without constituting a fully specified conjecture.

The source also asks what happens for ``intruded Aztec diamonds,'' and in particular why their tiling counts are so square-ish. The supplied text does not define this family, specify its relation to the intruded squares, or state an exact numerical assertion for it. Its wording presents the two-parameter assertion as a possibility, rather than an established fact.
\subsection{Short English statement}
Propp asks for an explanation of the square-like tiling counts of ``intruded Aztec diamonds'' and tentatively proposes that a square with an intrusion of length \(m\) always has a square or twice-square tiling count when its side length is \(2n\). The supplied excerpt omits the region construction, tiling rules, and parameter restrictions needed for a precise conjecture.
\subsection{Sources}
\begin{itemize}
\item Local UnsolvedMath registry, permanent ID \(800340\), source problem code \texttt{catalogue entry}, title ``Enumeration of Matchings --- Problem 11.'' The registry attributes the problem to Propp (1998) and labels it open; this is supplied status metadata, not an independently verified assessment.
\item Propp, \emph{Twenty Open Problems in Enumeration of Matchings} (1998), Problem 11 (\#11). Supplied source URL: \url{https://arxiv.org/abs/math/9801060}. The record reports access on 2026-10-16 and extraction by \texttt{arxiv\_latex}; these are supplied provenance details, not a claim of independent source access or verification.
\end{itemize}
\end{document}
